\documentclass[12pt]{article}
\usepackage{amsmath}
\usepackage{graphicx}
\usepackage{xcolor}
\definecolor{claudeorange}{HTML}{D97757}
\newenvironment{review}{\par\noindent\color{claudeorange}\textbf{Reviewer: }}{\par}
\newenvironment{addition-suggestion}{\par\noindent\color{claudeorange}\itshape\textbf{Suggested addition: }}{\par}

\title{Adsorption of Microplastics on Zeolite Frameworks: A Computational Study}
\author{J.A. Quiñonero}
\date{2026}

\begin{document}
\maketitle

\begin{abstract}
We present a computational study of the adsorption of polyethylene terephthalate (PET)
monomer units on ZSM-5 zeolite frameworks. Density functional theory calculations were
performed and adsorption energies were computed. The results show strong physisorption
at the channel intersection sites.
\end{abstract}

\section{Introduction}

Microplastic contamination in water systems is an emerging environmental concern.
Zeolites have been proposed as adsorbents due to their high surface area and tunable
pore chemistry. Previous work has examined simple organic molecules but not plastic
monomer units.

\section{Computational Methods}

All calculations were performed using the VASP software package.
The PBE functional was employed throughout. A plane-wave cutoff of 520 eV was used.
Dispersion corrections were added using the DFT-D3 scheme.

\section{Results}

Adsorption energies for the three adsorbates range from $-0.8$ to $-1.4$ eV.
The strongest binding is observed at the sinusoidal channel intersection.

\end{document}
